\chapter*{附录二 《圣祖仁皇帝实录》记载索额图获罪原因}\addcontentsline{toc}{chapter}{附录二 《圣祖仁皇帝实录》记载索额图获罪原因}\markboth{附录二 《圣祖仁皇帝实录》记载索额图获罪原因}{附录二 《圣祖仁皇帝实录》记载索额图获罪原因}
% 原书第378页

（康熙四十二年五月壬戌）领侍卫内大臣、和硕额驸尚之隆等传上谕：观索额图，并无退悔之意，背后怨尤，议论国事。伊之党类，朕皆访知。阿米达、麻尔图、额库礼、温待、邵甘、佟宝，伊等结党议论国事，威吓众人。


% 原书第378页

且索额图施威恐吓，举国之人尽惧索额图乎？亦有不惧者，即今索额图家人，已将伊告发。索额图能杀害乎？


% 原书第378页

至温待、额库礼，俱犯重罪流徙之人，因其年老，令回京师。伊等应安静以养余年，乃与索额图结党议论国事，妄自怨尤。伊等之党，俱属利口、愚昧无知之徒，被伊等恐吓，遂极畏惧。果至可杀之时，索额图能杀人或被人杀俱未可料。虽口称杀人，被杀者谁乎?


% 原书第378–379页

至索额图之党，汉官亦多。朕若尽指出，俱至灭族。朕不嗜杀人，嗣后尔等，若与索额图绝交，将所行之事举出，尚可。不然被索额图株连，必如噶褚哈、阿思哈被灭族之祸。索额图之党阿米达、额库礼、温待、麻尔图、佟宝、邵甘之同祖子孙、在部院者，俱查明革退。副都统佟宝不在家，俟到时再传谕此旨。着晓谕门上大人，与众侍卫等，尔等若在索额图处行走，必被索额图连累致死。


% 原书第379页

（癸亥）谕刑部，额库礼、温待、麻尔图、邵甘等，俱系问罪充发之犯。念其年老，取回京师。伊等至京，并不善养余年，与索额图结党、议论国事。俱着锁禁，不可疏放。阿米达姑念年老旧臣，着宽恕。


% 原书第379–380页

命近御侍卫海青，乾清门侍卫武格、马武传谕索额图：尔家人告尔之事，留内三年，朕有宽尔之意，尔并无退悔之意，背后仍怨尤，议论国事，结党妄行。尔背后怨尤之言，不可宣说，尔心内甚明。举国俱系受朕深恩之人，若受恩者半，不受恩者半，即俱从尔矣。去年皇太子在德州住时，尔乘马至皇太子中门方下，即此是尔应死处。尔自视为何等人，此处问公福善即知。尔任大学士时，因贪恶革退，后朕复起用，尔并不思念朕恩。即如养犬，尚知主恩，若尔者，极力加恩，亦属无益。朕差人搜与尔行走之江潢家，得书字甚多。朕亦欲差人到尔家搜看，但被尔连累之人甚多，举国俱不得安，所以中止。朕若不先发，尔必先之，朕亦熟思之矣。朕将尔行事指出一端，就可在此正法。尚念尔原系大臣，朕心不忍。但着尔闲住，又恐结党生事，背后怨尤议论。着交宗人府，与根度一处拘禁，不可疏放。


% 原书第380页

谕领侍卫内大臣等，索额图之子，并家内紧要人，俱交与心裕、发保，加意拘禁候旨。若其间别生事端，即将心裕、发保等族诛。

